\section*{[Prénom Nom]}

\subsection*{Purification et Witches' Brew : où passe le budget de l'attaque}

\paragraph{Idée.}
PureGen (Bhat et al., 2024) montre qu'on peut presque annuler Witches'
Brew (WB) en \emph{purifiant} les images d'entraînement avant d'apprendre,
c'est-à-dire en les faisant passer dans un filtre qui les ramène vers des
images naturelles (compression JPEG, débruitage par un modèle génératif).
Le succès de l'attaque tombe alors à presque 0\,\%. Pourtant, la seule
garantie théorique pour ce genre de défense (Hong et al., 2024) ne couvre
qu'une perturbation totale de norme $\ell_2$ d'environ $0{,}6$, alors que
WB peut dépenser jusqu'à $\sqrt{500}\times3{,}48\approx78$ sur CIFAR-10.
La théorie n'explique donc pas pourquoi la défense marche.

Mon hypothèse part de la façon dont une image se décompose en
fréquences. Les \emph{basses fréquences} portent les formes et les grandes
zones de couleur ; les \emph{hautes fréquences} portent les détails fins,
la texture et le bruit, pixel à pixel. Je pense que WB met une bonne
partie de sa perturbation dans les hautes fréquences, et que c'est
justement ce que la purification efface. Si c'est le cas, une version de
WB qui ne modifie que les basses fréquences devrait résister à la
purification.

\paragraph{Argument.}
Pour pouvoir faire le calcul, j'ai remplacé le purificateur par un filtre
simple, un \emph{filtre passe-bas}. J'utilise la transformée en cosinus
discrète (DCT, la même que dans JPEG), qui réécrit une image comme une
somme de motifs du plus lisse au plus fin. Je garde les $c\times c$ motifs
les plus lisses et je mets les autres à zéro (l'image fait
$224\times224$, donc $c=84$ revient à garder 14\,\% des motifs). Ce filtre
est une projection $P$ : la victime voit $Px+P\delta$ au lieu de
$x+\delta$, et la part de la perturbation qui arrive au modèle vaut
exactement $\|P\delta\|^2/\|\delta\|^2$ (théorème de Pythagore). Tout le
reste est effacé.

Pour un classifieur linéaire, le gradient de la perte est le produit d'un
vecteur d'erreur et de l'image : $\nabla_W\ell=(p-e_y)\,x^\top$, où $p$
est le vecteur des probabilités prédites et $e_y$ celui de la bonne
classe. Le cosinus que WB cherche à maximiser, entre le gradient de la
cible et celui d'un poison purifié $\tilde x=P(x+\delta)$, se coupe alors
en deux :
\[
\cos\big(\nabla_W\ell(x_t,y_{\mathrm{adv}}),\nabla_W\ell(\tilde x,y)\big)
=\underbrace{\cos(p_t-e_{y_{\mathrm{adv}}},\,\tilde p-e_y)}_{\text{les erreurs se ressemblent}}
\times\underbrace{\cos(x_t,\tilde x)}_{\text{les images se ressemblent}} .
\]
La perturbation n'agit sur ce cosinus qu'à travers $P\delta$ : tout ce
que WB a mis dans les hautes fréquences ne sert plus à rien une fois
l'image purifiée.

\paragraph{Ma variante de l'attaque : WB basses fréquences.}
Pour tester ça, j'ai modifié l'attaque. Après chaque pas d'optimisation,
je passe la perturbation $\delta$ dans le même filtre passe-bas, puis je
la ramène sous la limite de $16/255$ par pixel. Je répète ces deux étapes
quelques fois, car chacune peut défaire un peu l'autre. À la fin,
$\delta$ ne contient que des basses fréquences, donc le filtre ne lui
retire rien. J'appelle cette variante ``WB basses fréquences''.
Ma prédiction : sans défense, elle devrait être un peu moins efficace que
WB normal, parce qu'elle a moins de liberté ; avec la défense, elle
devrait être plus efficace.

\paragraph{Test 1 : ce qui survit à la purification (ImageNet).}
J'ai lancé WB normal et WB basses fréquences sur 3 cibles ImageNet
($224\times224$, ResNet-18 pré-entraîné, environ 1\,200 poisons). Pour
vérifier l'argument, j'ai mesuré la part de la perturbation qui survit au
filtre passe-bas, et le cosinus entre le gradient de la cible et celui des
poisons (la quantité que WB maximise), avant et après purification.
J'ai aussi testé une compression JPEG simplifiée (qualité 75 : l'image
est découpée en blocs de $8\times8$ pixels, chaque bloc passe en DCT, ses
coefficients sont arrondis selon la table de quantification standard de
JPEG, puis on revient en pixels). Les valeurs sont des moyennes sur les
3 cibles.

\begin{center}
\small
\begin{tabular}{@{}lcc@{}}
\toprule
& WB normal & WB basses fréq. \\
\midrule
Cosinus avant purification & 0{,}73 & 0{,}71 \\
Part de la perturbation restante après passe-bas & 0{,}34 & 1{,}00 \\
Cosinus après passe-bas & 0{,}38 & 0{,}61 \\
Alignement conservé après passe-bas & 53\,\% & 87\,\% \\
Alignement conservé après JPEG & 100\,\% & 99\,\% \\
\bottomrule
\end{tabular}
\end{center}

C'est cohérent avec la théorie. Après le passe-bas, WB normal ne garde
que 34\,\% de sa perturbation et la moitié de son alignement. WB basses
fréquences garde toute sa perturbation et 87\,\% de son alignement, alors
qu'il partait presque du même alignement : se limiter aux basses
fréquences coûte donc peu à l'attaque. JPEG ne change rien à cette
résolution, parce qu'il travaille sur des blocs de $8\times8$ pixels, ce
qui ne touche que des détails très fins dans une image de
$224\times224$.

\paragraph{Test 2 : taux de succès réel (Imagenette, budget 1\,\%).}
L'alignement ne dit pas si la cible bascule vraiment. J'ai donc croisé
deux choix : WB normal ou WB basses fréquences, sans défense ou avec le
filtre passe-bas ($c=84$). Ça donne 4 conditions. J'ai utilisé Imagenette
(un sous-ensemble d'ImageNet à 10 classes, environ 9\,500 images
d'entraînement), un ResNet-18 entraîné de zéro (40 époques, environ
82\,\% de précision) et les mêmes 10 cibles dans les 4 conditions. Je
regarde le succès et $L$, la perte finale de la cible vers la classe
visée par l'attaquant : plus $L$ est bas, plus la cible est proche de
basculer.

\begin{center}
\small
\begin{tabular}{@{}lcccc@{}}
\toprule
& WB normal & WB normal & WB basses fréq. & WB basses fréq. \\
& sans défense & + passe-bas & sans défense & + passe-bas \\
\midrule
Succès & 2/10 & 2/10 & 3/10 & 2/10 \\
$L$ moyen & 3{,}46 & 3{,}71 & 4{,}59 & 4{,}23 \\
Précision & 82\,\% & 80\,\% & 82\,\% & 81\,\% \\
\bottomrule
\end{tabular}
\end{center}

Les tendances vont dans le sens prévu : le passe-bas augmente $L$ pour WB
normal mais pas pour ma variante, et ma variante coûte un peu sans
défense. Mais rien n'est significatif : j'ai comparé les conditions
cible par cible avec un test de Wilcoxon, et toutes les différences ont
une probabilité d'au moins 20\,\% d'être dues au hasard. À 1\,\%,
l'attaque est trop faible : 7 cibles sur 10 échouent partout. Un bug de
graine m'a fait relancer les mêmes cibles une deuxième fois, ce qui m'a
servi de réplication : le résultat dépend surtout de la cible, et $L$
varie de 1 à 3 points entre deux runs identiques. J'ai aussi entraîné une
victime sans aucun poison : une des cibles est déjà classée dans la
classe visée sans attaque, donc un des succès ne venait pas de l'attaque.

\paragraph{Test 3 : budget 5\,\% et force de la défense.}
J'ai monté le budget à 5\,\% (474 poisons) et fait varier $c$, le nombre
de motifs gardés par le filtre, pour WB normal. Plus $c$ est petit, plus
le filtre efface. Les résultats sont partiels (1 à 4 cibles par ligne) :

\begin{center}
\small
\begin{tabular}{@{}lccc@{}}
\toprule
& Succès & Part de la perturbation effacée & Précision \\
\midrule
Sans défense & 1/2 & -- & 82\,\% \\
Passe-bas $c=84$ (garde 14\,\% des motifs) & 4/4 & 28 à 42\,\% & 82\,\% \\
Passe-bas $c=56$ (garde 6\,\% des motifs) & 1/2 & environ 50\,\% & 76 à 79\,\% \\
Passe-bas $c=28$ (garde 1{,}6\,\% des motifs) & 1/4 & environ 72\,\% & 60 à 66\,\% \\
\bottomrule
\end{tabular}
\end{center}

À 5\,\%, l'attaque devient beaucoup plus forte : une cible qui échouait
partout à 1\,\% bascule même avec le filtre $c=84$. Plus le filtre efface
la perturbation, plus WB échoue, ce qui correspond au mécanisme prévu.
Mais ce filtre est une défense trop grossière : il ne fait pas la
différence entre la perturbation et les détails utiles de l'image. À
$c=84$ il ne protège pas, et à $c=28$ il protège au prix de 20 points de
précision.

\paragraph{Où j'en suis.}
Les mesures du test 1 confirment la théorie. Sur le taux de succès réel,
je n'ai pas encore de résultat concluant. J'ai lancé les 4 conditions à
5\,\% avec $c=56$, sur 10 cibles par condition.
